\hypertarget{chapter-13-indexes-glossaries-and-nomenclature}{%
\chapter{Indexes, Glossaries, and Nomenclature}\label{chapter-13-indexes-glossaries-and-nomenclature}}

A book's index and glossary are the parts of its back matter readers consult rather than read straight through, and precisely because of that, they are easy to treat as an afterthought bolted on once the main text is finished. For a project with its own recurring, specific vocabulary, built up across many chapters or many books, treating index and glossary construction as an afterthought produces exactly the kind of inconsistency that makes them least useful: the same concept indexed under two different terms, or a glossary that defines a term without covering every sense the book actually used it in.

\hypertarget{building-an-index-with-makeindex}{%
\section{\texorpdfstring{Building an Index with \texttt{makeindex}}{Building an Index with makeindex}}\label{building-an-index-with-makeindex}}

An index entry in LaTeX is marked at the point in the text it applies to, using \texttt{\textbackslash{}index\{term\}}, and the actual sorted, page-numbered index is generated by a separate tool, \texttt{makeindex}, run on the raw index data LaTeX writes out during compilation, in the same multi-pass pattern Chapter 5 described for bibliographies: LaTeX writes raw entries, \texttt{makeindex} processes and sorts them, and a subsequent LaTeX pass reads the processed result back in and typesets it.

\texttt{\textbackslash{}index} entries support subentries (\texttt{\textbackslash{}index\{notation!custom\ operators\}}, producing a nested index entry under ``notation''), cross-references (\texttt{\textbackslash{}index\{notation\textbar{}see\{symbols\}\}}), and explicit sort keys for terms that should be alphabetized differently from how they are typeset (\texttt{\textbackslash{}index\{TeX@\textbackslash{}TeX\}}, ensuring a command like \texttt{\textbackslash{}TeX} sorts under T rather than by its literal, unsortable macro name). For a technical book with its own recurring vocabulary, subentries are worth using deliberately rather than only when a term happens to accumulate enough page references to need them: grouping related terms under a shared parent entry from the start makes the index itself into a small map of the book's conceptual structure, not merely an alphabetized list of page numbers.

The discipline that actually determines whether an index is useful is consistency at the point of marking, not anything \texttt{makeindex} itself can enforce. A concept indexed as ``admissibility'' in one chapter and ``admissible states'' in another produces two separate index entries for what a reader experiences as one concept, and no tool catches this automatically; it has to be caught by maintaining a short list of the book's key indexed terms, exactly as Chapter 10 recommended maintaining a shared notation file, and checking new \texttt{\textbackslash{}index} entries against that list rather than inventing a phrasing fresh at each point in the text.

\hypertarget{glossaries-and-nomenclature-with-the-glossaries-package}{%
\section{\texorpdfstring{Glossaries and Nomenclature with the \texttt{glossaries} Package}{Glossaries and Nomenclature with the glossaries Package}}\label{glossaries-and-nomenclature-with-the-glossaries-package}}

The \texttt{glossaries} package handles a related but distinct need: defining a term once, with a full definition, and then referencing it throughout the text in a way that can produce both an inline reference and a collected glossary at the back of the book. A term is declared once, typically in the shared preamble file described in Chapter 3, with \texttt{\textbackslash{}newglossaryentry\{admissibility\}\{name=\{admissibility\},\ description=\{...\}\}}, and referenced in the body with \texttt{\textbackslash{}gls\{admissibility\}}, which typesets the term itself while also recording that this location used it, feeding into the automatically generated glossary the same way indexed terms feed into the index.

\texttt{glossaries} additionally supports first-use behavior distinct from subsequent uses, useful for acronyms in particular: \texttt{\textbackslash{}gls} on an acronym's first use in the document can expand to the full term followed by the abbreviation in parentheses, while every subsequent use renders as the abbreviation alone, without the author needing to track by hand which use is the first. For a book defining a genuinely technical vocabulary rather than only abbreviations, the same first-use mechanism can be used more broadly, so a term's first appearance in each chapter, or in the book as a whole depending on configuration, gets slightly more context than later uses that can now rely on the reader already having encountered the definition.

A nomenclature list, distinct from a glossary of terms, is more appropriate for a book with heavy recurring mathematical notation: rather than defining what a word means, it lists what a symbol means, typically organized by the order symbols were introduced or grouped by category (Latin letters, Greek letters, custom operators). The \texttt{nomencl} package, or \texttt{glossaries}'s own support for symbol-type entries, both handle this, and for a project already maintaining the shared notation file recommended in Chapter 11, generating a nomenclature list from that same file, rather than maintaining the list of symbols separately, keeps the nomenclature from drifting out of sync with the notation actually used in the text.

\hypertarget{when-a-glossary-earns-its-keep}{%
\section{When a Glossary Earns Its Keep}\label{when-a-glossary-earns-its-keep}}

Not every book needs a glossary, and adding one reflexively to any technical project produces the same kind of clutter an unnecessary index would: a glossary of five terms a reader could infer from context on first encounter adds a page of back matter without adding proportional value, while a genuinely large, recurring technical vocabulary, the kind a reader is likely to need to look back up multiple times across a long book, is exactly the case a glossary is built for.

A reasonable test, worth applying deliberately rather than defaulting either way, is whether a term is likely to recur across chapters written far enough apart, in drafting time, that a reader (or the author, months later) would plausibly want to check its definition again without rereading the chapter that first introduced it. Terms that clear this bar belong in a glossary; terms defined and used entirely within a single chapter, never referenced again, are usually better served by a definition environment at the point of use, as covered in Chapter 11, than by a glossary entry that would only ever be consulted from the one place it is used.

\hypertarget{nomenclature-across-a-multi-book-corpus}{%
\section{Nomenclature Across a Multi-Book Corpus}\label{nomenclature-across-a-multi-book-corpus}}

For a research program spanning many books that share notation, as Chapter 10 and Chapter 11 both recommended treating symbol definitions as a maintained cross-corpus resource, the nomenclature list built for any one book benefits from being generated from that same shared resource rather than compiled independently per book. A symbol's meaning that stays consistent across the corpus should produce an identical nomenclature entry wherever it appears, and generating that entry automatically from the single shared notation file, rather than transcribing it by hand into each new book's own nomenclature list, is both less work and the more reliable way to keep a large body of related books actually consistent with each other over years of continued writing.

Part V closes here, with the reference apparatus, bibliography, index, glossary, and nomenclature, all treated as maintained resources rather than one-off additions. Part VI turns to graphics and automation: bringing figures, generated plots, and computation itself directly into the publishing pipeline this book has been building toward since Chapter 3.
